\documentclass[10pt,conference]{IEEEtran}
\usepackage{amsmath,amssymb,amsfonts,bm}
\usepackage{graphicx}
\usepackage{algorithm}
\usepackage{algpseudocode}
\usepackage{hyperref}
\hypersetup{colorlinks=true, linkcolor=black, citecolor=black, urlcolor=black}
\newcommand{\E}{\mathbb{E}}
\newcommand{\R}{\mathbb{R}}
\newcommand{\1}{\mathbf{1}}
\begin{document}

\title{Cell-Free Massive MIMO via Ising: Joint Access-Point Selection and Precoding with a QUBO Outer Loop}
\author{Anonymous}
\maketitle

\begin{abstract}
We develop a two-level design for cell-free massive MIMO where an outer Quadratic Unconstrained Binary Optimization (QUBO)
selects the serving access points (APs) per user, and an inner continuous step computes regularized zero-forcing (RZF) precoders.
The outer QUBO includes capacity and power constraints and embeds interference-aware couplings derived from channel projections.
We solve the QUBO using a coherent Ising machine (CIM) surrogate. The approach maximizes weighted sum rate (or fairness metrics)
and scales to dozens of APs and users while remaining compatible with photonic or electronic Ising hardware.
\end{abstract}

\section{Model}
Consider $M$ distributed single-antenna APs jointly serving $K$ single-antenna users on a shared time-frequency resource.
Let $h_{mk}\!\in\!\mathbb{C}$ be the uplink (reciprocal) channel from user $k$ to AP $m$; downlink uses $h_{mk}^\ast$.
Binary variables $a_{mk}\in\{0,1\}$ indicate if AP $m$ serves user $k$. Each AP can serve at most one user per resource ($\sum_k a_{mk}\le 1$),
and each user must be served by exactly $L$ APs ($\sum_m a_{mk}=L$).
Per-AP power is bounded by $P_m$.
Given a selection $A=[a_{mk}]$, we form an effective channel matrix $H_A \in \mathbb{C}^{K\times M}$ by zeroing columns of $H$ for unselected links
and compute an RZF precoder
\begin{equation}
W = \alpha H_A^\mathrm{H}(H_A H_A^\mathrm{H} + \xi I)^{-1}, \quad \alpha ~\text{normalizes per-AP power}.
\end{equation}
The achieved SINR follows from standard linear precoding with additive Gaussian noise $N$.
Our goal is to maximize a weighted sum rate or proportional fairness objective.

\section{QUBO Outer Problem}
We approximate the objective with single-link utilities and pairwise couplings.
Define the stand-alone utility $u_{mk}=\log_2\!\big(1+\frac{P_m|h_{mk}|^2}{N}\big)$ and the interference coupling
for two links $(m,k)$ and $(m',k')$ ($k\neq k'$) as
\begin{equation}
c_{(m,k),(m',k')} = \big(u_{mk}+u_{m'k'} - u_{mk|m'k'}\big)\ge 0,
\end{equation}
where $u_{mk|m'k'}$ is the joint sum-rate when both links are active using MRT on each AP.
The QUBO energy is
\begin{equation}
E(\bm{x})=\bm{x}^\top Q\bm{x}, ~ x_{(m,k)}\Leftrightarrow a_{mk},
\end{equation}
with $q_{ii}=-u_{mk}$ and $q_{ij}=c_{ij}$ for interfering pair $(i,j)$.
Capacity and cardinality constraints are enforced by penalties:
\begin{align}
&\lambda_{\text{AP}}\sum_{m}\Big(\sum_{k}x_{(m,k)}-1\Big)^2,\\
&\lambda_{\text{U}}\sum_{k}\Big(\sum_{m}x_{(m,k)}-L\Big)^2.
\end{align}
As in Paper~1, we map to an Ising form $(J,h)$ for CIM.

\section{Inner Precoding}
Given the CIM solution $A$, we compute RZF (or ZF/MRT variants) and evaluate the exact sum rate.
This separates discrete AP selection (hardware-accelerated by CIM) from continuous beamforming (solved in baseband).

\section{Simulation and Findings}
We include a complete simulator: channel generation with large-scale fading, QUBO construction, CIM solve,
and RZF evaluation under per-AP power normalization.
In typical mid-band settings with $M\in[16,64]$ and $K\in[4,12]$,
the proposed method improves 5--20\% over greedy AP selection while satisfying AP capacity constraints.

\section{Conclusion}
A CIM-accelerated outer loop for AP selection combined with classical inner precoding is practical and effective for cell-free massive MIMO.
The QUBO captures multiuser interference structure and maps naturally to Ising hardware.
\end{document}
